\begin{lemma}[Stability of the limiting equation under locally uniform limits]
Let $F$ be a finite-dimensional real inner-product space, let $\Omega$ be an
open subset of the space-time $\ell^2(F\times\R)$, let $v_1,v_2,\ldots$ each be a
viscosity solution of $F[v_j]=0$ on $\Omega$ in the sense of
\noderef{IsLimitSolutionOn}, and let $v_0$ be a real function on space-time such
that $v_j\to v_0$ uniformly on every compact subset of $\Omega$. Then $v_0$ is a
viscosity solution of $F[v_0]=0$ on $\Omega$.
\end{lemma}
\begin{proof}
\emph{Step 0: continuity of the limit.}
Let $\zeta\in\Omega$. Since $\Omega$ is open there is $\varrho>0$ with
$\overline B_\varrho(\zeta)\subseteq\Omega$, a compact set, so $v_j\to v_0$
uniformly on it. By \noderef{IsLimitSolutionOn} each $v_j$ is continuous on
$\Omega$, hence on $\overline B_\varrho(\zeta)$, and a uniform limit of
continuous functions on a set is continuous on that set. Thus $v_0$ is
continuous at every point of $\Omega$, i.e. continuous on $\Omega$.

\emph{Step 1: strictifying an upper test.}
Let $\zeta_0\in\Omega$ and let $\phi$ be $C^2$ near $\zeta_0$ with
$\zeta\mapsto v_0(\zeta)-\phi(\zeta)$ having a local maximum on $\Omega$ at
$\zeta_0$. Set
\[
  \phi^\sharp(\zeta)=\phi(\zeta)+|\zeta-\zeta_0|^4 .
\]
The function $\zeta\mapsto|\zeta-\zeta_0|^4=\bigl(|\zeta-\zeta_0|^2\bigr)^2$ is a
polynomial in the coordinates, hence smooth, and
\begin{equation}\label{eq:ls-fourth}
  \nabla\bigl(|\cdot-\zeta_0|^4\bigr)(\zeta_0)=0,
  \qquad
  D^2\bigl(|\cdot-\zeta_0|^4\bigr)(\zeta_0)=0 ,
\end{equation}
since its gradient is $4|\zeta-\zeta_0|^2(\zeta-\zeta_0)$, which vanishes at
$\zeta_0$ together with its derivative. Hence $\phi^\sharp$ is $C^2$ near
$\zeta_0$ with $\nabla\phi^\sharp(\zeta_0)=\nabla\phi(\zeta_0)$ and
$D^2\phi^\sharp(\zeta_0)=D^2\phi(\zeta_0)$, so by \noderef{LimitOp}
\begin{equation}\label{eq:ls-same-jet}
  F[\phi^\sharp](\zeta_0)=F[\phi](\zeta_0).
\end{equation}

Choose $\varrho>0$ so small that $\overline B_\varrho(\zeta_0)\subseteq\Omega$,
that $\phi$ is $C^2$ on a neighbourhood of $\overline B_\varrho(\zeta_0)$, and
that $v_0-\phi$ attains its maximum over $\overline B_\varrho(\zeta_0)$ at
$\zeta_0$; the last is possible because the maximum is local on $\Omega$. Then
for $\zeta\in\overline B_\varrho(\zeta_0)$ with $\zeta\ne\zeta_0$,
\[
  v_0(\zeta)-\phi^\sharp(\zeta)
   =\bigl(v_0(\zeta)-\phi(\zeta)\bigr)-|\zeta-\zeta_0|^4
   <v_0(\zeta_0)-\phi(\zeta_0)=v_0(\zeta_0)-\phi^\sharp(\zeta_0),
\]
so $\zeta_0$ is the \emph{unique} maximum point of $v_0-\phi^\sharp$ on
$\overline B_\varrho(\zeta_0)$.

\emph{Step 2: nearby maxima of the approximations.}
The set $\overline B_\varrho(\zeta_0)$ is compact and contained in $\Omega$, and
each $v_j-\phi^\sharp$ is continuous on it, so it attains a maximum at some
$\zeta_j\in\overline B_\varrho(\zeta_0)$. We claim $\zeta_j\to\zeta_0$. By
compactness it suffices to show that every convergent subsequence has limit
$\zeta_0$. Let $\zeta_{j_k}\to\bar\zeta\in\overline B_\varrho(\zeta_0)$. Uniform
convergence of $v_j$ to $v_0$ on $\overline B_\varrho(\zeta_0)$ and continuity of
$v_0-\phi^\sharp$ there give
\[
  (v_{j_k}-\phi^\sharp)(\zeta_{j_k})\longrightarrow
  (v_0-\phi^\sharp)(\bar\zeta).
\]
On the other hand, for every $\zeta\in\overline B_\varrho(\zeta_0)$,
$(v_{j_k}-\phi^\sharp)(\zeta_{j_k})\ge(v_{j_k}-\phi^\sharp)(\zeta)
 \to(v_0-\phi^\sharp)(\zeta)$, so
$(v_0-\phi^\sharp)(\bar\zeta)\ge(v_0-\phi^\sharp)(\zeta)$ for all such $\zeta$.
Thus $\bar\zeta$ is a maximum point of $v_0-\phi^\sharp$ on
$\overline B_\varrho(\zeta_0)$, and by the uniqueness in Step 1,
$\bar\zeta=\zeta_0$. Hence $\zeta_j\to\zeta_0$; in particular
$|\zeta_j-\zeta_0|<\varrho$ for all large $j$, so $\zeta_j$ is an interior point
of $\overline B_\varrho(\zeta_0)$ and therefore $v_j-\phi^\sharp$ has a local
maximum on $\Omega$ at $\zeta_j$.

\emph{Step 3: passing the inequality.}
For such $j$, $\phi^\sharp$ is $C^2$ near $\zeta_j$ and $v_j-\phi^\sharp$ has a
local maximum on $\Omega$ at $\zeta_j$, so the subsolution inequality of
\noderef{IsLimitSubOn}, available because $v_j$ is a solution by
\noderef{IsLimitSolutionOn}, gives
\[
  F[\phi^\sharp](\zeta_j)\ge0 .
\]
Since $\phi^\sharp$ is $C^2$ on a neighbourhood of
$\overline B_\varrho(\zeta_0)$, its gradient and second derivative are
continuous there; by \noderef{LimitOp} the map
$\zeta\mapsto F[\phi^\sharp](\zeta)$, which is the second derivative of
$\phi^\sharp$ at $\zeta$ evaluated on the pair $(\xi_\zeta,\xi_\zeta)$ with
$\xi_\zeta=(\nabla_x\phi^\sharp(\zeta),1)$, is therefore continuous near
$\zeta_0$. Letting $j\to\infty$ and using $\zeta_j\to\zeta_0$ gives
$F[\phi^\sharp](\zeta_0)\ge0$, and \eqref{eq:ls-same-jet} yields
$F[\phi](\zeta_0)\ge0$.

\emph{Step 4: lower tests.}
If instead $v_0-\phi$ has a local minimum on $\Omega$ at $\zeta_0$, repeat
Steps 1--3 with $\phi^\sharp(\zeta)=\phi(\zeta)-|\zeta-\zeta_0|^4$, with minima
in place of maxima, and with the supersolution inequality of
\noderef{IsLimitSuperOn}. Identity \eqref{eq:ls-fourth} again gives
$F[\phi^\sharp](\zeta_0)=F[\phi](\zeta_0)$, and the conclusion is
$F[\phi](\zeta_0)\le0$.

Together with Step 0 and the openness of $\Omega$, this verifies both clauses of
\noderef{IsLimitSolutionOn} for $v_0$.
\end{proof}
